\documentclass[10pt,twocolumn]{article}
\usepackage[margin=1in]{geometry}
\usepackage{graphicx}
\usepackage{booktabs}
\usepackage{amsmath}
\usepackage{hyperref}
\usepackage{caption}
\usepackage{subcaption}
\usepackage{float}
\usepackage{cite}
\usepackage{microtype}
\usepackage{parskip}
\usepackage{titlesec}
\titlespacing*{\section}{0pt}{6pt}{3pt}
\titlespacing*{\subsection}{0pt}{4pt}{2pt}

\graphicspath{{results/}}

% Compact float spacing to reduce large gaps around figures and tables
\captionsetup{font=small,skip=3pt}
\setlength{\textfloatsep}{6pt plus 1pt minus 2pt}
\setlength{\floatsep}{6pt plus 1pt minus 2pt}
\setlength{\intextsep}{6pt plus 1pt minus 2pt}
\setlength{\dbltextfloatsep}{6pt plus 1pt minus 2pt}
\raggedbottom

\title{\textbf{Adaptive Local Search Operator Selection for the\\
Traveling Salesperson Problem via Tabular Q-Learning \\AIE635 Project}}
\author{\textit{Akyaz Ozlem, Erdas Sila}}
\date{}

\begin{document}
\maketitle

%--------------------------------------------------------------------
\begin{abstract}
The Traveling Salesperson Problem (TSP) is a canonical NP-hard combinatorial
optimization problem. Local search operators such as 2-opt, swap, and relocate
are effective in isolation, yet no single operator dominates across all phases
of a search trajectory. This paper proposes a \emph{selection hyper-heuristic}
that frames operator selection as a Markov Decision Process (MDP) and learns a
dynamic policy through tabular Q-learning. Our 120-state formulation encodes
solution quality, search phase, recency of improvement, and last operator used.
To ensure mathematical rigor, all evaluations strictly adhere to the TSPLIB EUC\_2D
distance standard. We compare three $\varepsilon$-greedy exploration strategies 
and two reward definitions on four standard TSPLIB instances over 30 independent runs. 
The Q-learning agent achieves a near-optimal mean gap across all instances (e.g., 3.87\% on st70 and 6.66\% on kroA100), outperforming all fixed-operator and random baselines. A novel
transition-matrix analysis reveals a powerful, universally autonomous operator synergy:
\emph{swap}$\to$\emph{2-opt} across all instances, successfully mimicking expert-designed
heuristics like Lin-Kernighan.
\end{abstract}

%--------------------------------------------------------------------
\section{Introduction}

The Traveling Salesperson Problem asks for the shortest Hamiltonian cycle
through $n$ cities. Despite its simple formulation, TSP is NP-hard~\cite{garey79},
and exact methods become computationally infeasible beyond a few hundred cities.
Local search heuristics remain the backbone of state-of-the-art solvers such as
Concorde and LKH~\cite{applegate06,helsgott00}.

A fundamental limitation of classical local search is the static use of a
single operator. In practice, operator effectiveness is highly context-sensitive:
2-opt efficiently untangles edges near optimal solutions, while perturbation
operators (swap, relocate) are needed to escape deep local optima. Manually
scheduling operators is time-consuming and instance-specific.

\emph{Selection hyper-heuristics}~\cite{burke10} address this by maintaining a
high-level policy that chooses among low-level heuristics at each step. When
the policy is \emph{learned}, the framework adapts online to the evolving
search landscape. Q-learning~\cite{watkins92} provides a natural mechanism:
the agent observes state, selects an operator, receives an improvement reward,
and updates a value function guiding future decisions.

This paper makes four contributions:
\begin{enumerate}\itemsep0pt
  \item A complete tabular Q-learning hyper-heuristic with a 120-state MDP.
  \item Systematic comparison of three exploration strategies and two reward definitions over 30 robust runs.
  \item A \textbf{transition-matrix analysis} quantifying autonomous operator synergies.
  \item A generalization study evaluating policy transfer across instance sizes, identifying oracle-state limitations.
\end{enumerate}

%--------------------------------------------------------------------
\section{Related Work}

\textbf{Local Search for TSP.}
The 2-opt heuristic~\cite{lin65} remains the workhorse of TSP local search.
Its extension to the Lin-Kernighan (LK) move~\cite{lin73} yields solutions
within 1--2\% of optimal on large instances. Or-opt (relocate) and city-swap
operators complement 2-opt by enabling moves it cannot perform~\cite{or76}.

\textbf{Hyper-heuristics.}
Burke et al.~\cite{burke10} provide a comprehensive taxonomy distinguishing
\emph{selection} from \emph{generation} hyper-heuristics. Our work falls in
the learning selection category. Meignan et al.~\cite{meignan15} apply
Q-learning to neighbourhood selection in vehicle routing; Da Costa et
al.~\cite{dacosta08} use multi-armed bandits for evolutionary operator
selection. Compared to these, we introduce a richer state encoding and a
transition-matrix diagnostic tool.

\textbf{Deep RL for Combinatorial Optimization.}
Bello et al.~\cite{bello16} (pointer networks) and Kool et al.~\cite{kool19}
(attention model) demonstrate that deep RL can construct near-LKH solutions,
but require GPU-scale training. Our tabular Q-learning is deliberately
lightweight and interpretable.

\textbf{Adaptive $\varepsilon$-Greedy.}
Tokic~\cite{tokic10} proposes adaptive epsilon based on value differences in
robot navigation. We adapt this idea to TSP operator selection via
stagnation detection.

%--------------------------------------------------------------------
\section{Methodology}

\subsection{Distance Metric and MDP Formulation}

To prevent measurement illusions caused by standard floating-point rounding, our 
methodology strictly utilizes the exact TSPLIB \texttt{EUC\_2D} distance standard 
($\lfloor \text{dist} + 0.5 \rfloor$), ensuring mathematically sound gap calculations.

We model the search as an episodic MDP $(\mathcal{S}, \mathcal{A}, \mathcal{R},
\mathcal{T}, \gamma)$ with $|\mathcal{S}|=120$ states.

Each state is a 4-tuple $(\text{quality\_band},\, \text{improved},\,
\text{stage},\, \text{last\_op})$:
\begin{itemize}\itemsep0pt
  \item \textbf{quality\_band} $\in\{0,..,4\}$: relative gap to optimal
        discretized as $<$5\%, 5--15\%, 15--30\%, 30--50\%, $>$50\%.
  \item \textbf{improved} $\in\{0,1\}$: did the last step improve the tour?
  \item \textbf{stage} $\in\{0,1,2\}$: Early ($<$33\%), Mid, Late search phase.
  \item \textbf{last\_op} $\in\{0,1,2,3\}$: previously applied operator.
\end{itemize}

Action space: $\mathcal{A}=\{$2-opt, swap, relocate, restart$\}$.

Two reward functions are compared:
\begin{align}
  R_1 &= \ell_{\text{prev}} - \ell_{\text{new}} \\
  R_2 &= R_1 + 100\cdot\mathbf{1}[\ell_{\text{new}} < \ell_{\text{best}}]
\end{align}

The Q-table is updated using the standard off-policy Bellman equation:
\begin{equation}
Q(s,a) \leftarrow Q(s,a) + \alpha\bigl[r + \gamma\max_{a'}Q(s',a') - Q(s,a)\bigr]
\end{equation}
with $\alpha=0.1$, $\gamma=0.9$.

\subsection{Operators}

\begin{table}[!htbp]
\centering
\caption{Local Search Operators}
\small
\begin{tabular}{ll}
\toprule
\textbf{Operator} & \textbf{Description} \\
\midrule
2-opt & Stochastic best-improvement reversal \\
Swap & Exchange two random cities \\
Relocate & Remove city; reinsert at different position \\
Restart & Full random permutation (perturbation) \\
\bottomrule
\end{tabular}
\end{table}

To maintain computational speed, \emph{2-opt} utilizes a stochastic best-improvement 
approach by sampling up to 100 random segment pairs, rather than performing an exhaustive 
$\mathcal{O}(N^2)$ search. Improving moves are always accepted; worsening moves are accepted with
probability 0.01 (SA-style escape), \emph{except} after restart which
unconditionally replaces the current tour.

\subsection{Exploration Strategies}

\begin{itemize}\itemsep0pt
  \item \textbf{Constant} $\varepsilon=0.1$: fixed throughout.
  \item \textbf{Exponential Decay}: $\varepsilon_0{=}1.0$, $\varepsilon_{\min}{=}0.01$,
        decay$=0.998$.
  \item \textbf{Adaptive Stagnation}: decays by $\times0.99$ after improvements;
        increases by $+0.05$ after 50 consecutive non-improving steps
        (capped at $\varepsilon_{\max}{=}0.8$).
\end{itemize}

\subsection{Transition Matrix Analysis}

After each episode we record $C[i,j]$: the number of times operator $j$ was
applied immediately after a \emph{successful} (improving) application of
operator $i$, excluding restart. Row-normalizing $C$ yields transition
probabilities revealing operator synergies.

\subsection{Experimental Setup}

\begin{table}[!htbp]
\centering
\caption{Experimental Parameters}
\small
\begin{tabular}{ll}
\toprule
\textbf{Parameter} & \textbf{Value} \\
\midrule
Instances & eil51 (opt=426), berlin52 (opt=7542) \\
          & st70 (opt=675), kroA100 (opt=21282) \\
Max iterations / episode & 3,000 \\
Runs per configuration & 30 \\
Training instances & eil51, berlin52 \\
Test instances & st70, kroA100 \\
\bottomrule
\end{tabular}
\end{table}

%--------------------------------------------------------------------
\section{Results and Discussion}

\subsection{Main Performance Comparison}

Table~\ref{tab:main} reports the average optimality gap (\%) over 30 independent runs,
ensuring statistical significance. Q-Learning comprehensively dominates across the instances. 
The agent's capacity to contextually shift operators enables it to vastly outperform 
both random operator selection and static fixed-operator baselines, reducing the gap on 
st70 and kroA100 to 3.87\% and 6.66\% respectively.


\begin{table}[!t]
\centering
\caption{Average Gap to Optimal (\%) $\pm$ Std. (30 Runs)}
\label{tab:main}
\scriptsize
\setlength{\tabcolsep}{2pt}
\renewcommand{\arraystretch}{0.92}
\resizebox{\columnwidth}{!}{%
\begin{tabular}{lcccc}
\toprule
\textbf{Strategy} & \textbf{eil51} & \textbf{berlin52} & \textbf{st70} & \textbf{kroA100} \\
\midrule
Q-Learning (Adap.) & \textbf{2.14}$\pm$0.82 & \textbf{0.34}$\pm$0.95 & \textbf{3.87}$\pm$1.18 & \textbf{6.66}$\pm$1.69 \\
Fixed 2-opt        & 3.94$\pm$1.55 & 2.17$\pm$2.41 & 5.42$\pm$1.80 & 7.10$\pm$0.62 \\
Fixed swap         & 22.35$\pm$3.19 & 17.12$\pm$4.17 & 27.70$\pm$5.31 & 27.72$\pm$3.59 \\
Fixed relocate     & 19.62$\pm$2.55 & 17.20$\pm$4.10 & 27.08$\pm$5.80 & 27.51$\pm$4.88 \\
Fixed restart      & 24.04$\pm$4.25 & 17.50$\pm$4.30 & 27.97$\pm$5.51 & 28.27$\pm$4.07 \\
Random select.     & 4.23$\pm$1.48 & 1.82$\pm$2.26 & 5.36$\pm$2.01 & 12.30$\pm$3.10 \\
\bottomrule
\end{tabular}%
}
\end{table}

Figure~\ref{fig:convergence} shows convergence curves. The Q-Learning algorithm rapidly 
descends below all baselines. On larger instances, the adaptive policy successfully avoids 
the deep stagnation regions where static methods get trapped.


\begin{figure*}[!t]
  \centering
  \begin{subfigure}{0.48\textwidth}
    \includegraphics[width=\linewidth]{convergence_eil51.png}
    \caption{eil51}
  \end{subfigure}
  \hfill
  \begin{subfigure}{0.48\textwidth}
    \includegraphics[width=\linewidth]{convergence_berlin52.png}
    \caption{berlin52}
  \end{subfigure}

  \vspace{2pt}
  \begin{subfigure}{0.48\textwidth}
    \includegraphics[width=\linewidth]{convergence_st70.png}
    \caption{st70}
  \end{subfigure}
  \hfill
  \begin{subfigure}{0.48\textwidth}
    \includegraphics[width=\linewidth]{convergence_kroA100.png}
    \caption{kroA100}
  \end{subfigure}
  \caption{Convergence curves: best tour length vs.\ iteration. Red dashed line = known optimal.}
  \label{fig:convergence}
\end{figure*}

\subsection{Autonomous Operator Synergy Discovery}

One of the most remarkable findings derived from the transition probability heatmaps 
(Figure~\ref{fig:heatmap}) is the agent's autonomous discovery of complex heuristic sequences. 
Across all evaluated instances, the Q-agent learned an overwhelmingly strong synergy: 
\emph{swap}$\to$\emph{2-opt} (prob.\ 1.00).

This signifies that the agent organically learned a foundational principle of advanced TSP 
metaheuristics (such as Lin-Kernighan): utilizing a perturbation operator (\emph{swap}) to 
dislodge the tour from a local optimum, immediately followed by an aggressive refinement 
operator (\emph{2-opt}) to optimize the newly exposed edges. The agent achieved this purely 
through reward-driven exploration without human injection of domain knowledge.


\begin{figure*}[!t]
  \centering
  \begin{subfigure}{0.48\textwidth}
    \includegraphics[width=\linewidth]{transition_heatmap_eil51.png}
    \caption{eil51}
  \end{subfigure}
  \hfill
  \begin{subfigure}{0.48\textwidth}
    \includegraphics[width=\linewidth]{transition_heatmap_kroA100.png}
    \caption{kroA100}
  \end{subfigure}
  \caption{Operator transition probability heatmaps after successful moves (restart excluded). Darker cells indicate stronger synergy.}
  \label{fig:heatmap}
\end{figure*}

\subsection{Reward Shaping: $R_1$ vs.\ $R_2$}

Both reward functions yielded \emph{identical} results. When the agent is far from the optimum, 
nearly every improving step also sets a new best, making $R_1$ and $R_2$ equivalent.
Near convergence, the +100 bonus is too infrequent to shift Q-values.

\begin{figure}[!htbp]
  \centering
  \includegraphics[width=0.72\linewidth]{reward_bar_comparison.png}
  \caption{Average gap for $R_1$ (Pure) vs.\ $R_2$ (Global Bonus) on
           training instances. Error bars = 1 std.}
  \label{fig:rewardbar}
\end{figure}

\begin{table}[!htbp]
\centering
\caption{Epsilon Strategy Results}
\label{tab:eps}

\footnotesize
\setlength{\tabcolsep}{3pt}

\begin{tabular}{lcccc}
\toprule
\textbf{Strategy} & \multicolumn{2}{c}{\textbf{Gap (\%)}} & \multicolumn{2}{c}{\textbf{Restarts}} \\
 & eil51 & berlin52 & eil51 & berlin52 \\
\midrule
Constant $\varepsilon{=}0.1$ & 3.47$\pm$1.15 & 2.04$\pm$1.05 & 74.8 & 69.4 \\
Exp. Decay                   & 4.84$\pm$0.87 & 3.19$\pm$1.32 & 134.2 & 122.0 \\
\textbf{Adaptive (Stag.)}    & \textbf{2.14}$\pm$0.82 & \textbf{0.34}$\pm$0.95 & 116.2 & 110.2 \\
\bottomrule
\end{tabular}
\end{table}

Exponential Decay performs worst despite starting at $\varepsilon{=}1.0$.
As shown in Figure~\ref{fig:epsiloncomp}, epsilon collapses to near-zero
within $\approx$1,000 iterations, causing the agent to exploit an immature
Q-table for the remaining budget. The high restart count confirms
repeated escape attempts from early-locked suboptimal policies.

Adaptive Stagnation maintains a meaningful $\varepsilon$ throughout, increasing
exploration precisely when stagnation is detected (sawtooth pattern visible in
Figure~\ref{fig:epsiloncomp}, right panel), and decaying after improvements.

\begin{figure}[!htbp]
  \centering
  \includegraphics[width=0.86\linewidth]{epsilon_comparison_eil51.png}
  \caption{Epsilon strategy comparison on eil51. Left: convergence curves.
           Right: $\varepsilon$ over time showing Adaptive's sawtooth pattern.}
  \label{fig:epsiloncomp}
\end{figure}

\begin{figure}[!htbp]
  \centering
  \includegraphics[width=0.78\linewidth]{epsilon_restart_bar.png}
  \caption{Average restart count per epsilon strategy. Exponential Decay
           triggers the most restarts, indicating policy instability.}
  \label{fig:restartbar}
\end{figure}

\subsection{Q-Table Analysis}

\begin{figure}[!htbp]
  \centering
  \begin{subfigure}{0.47\linewidth}
    \includegraphics[width=\linewidth]{qtable_heatmap_eil51.png}
    \caption{eil51}
  \end{subfigure}
  \hfill
  \begin{subfigure}{0.47\linewidth}
    \includegraphics[width=\linewidth]{qtable_heatmap_berlin52.png}
    \caption{berlin52}
  \end{subfigure}
  \caption{Q-values aggregated by solution quality band. Green = high Q-value
           (preferred operator). Darker red = avoided.}
  \label{fig:qtable}
\end{figure}

The Q-table heatmaps (Figure~\ref{fig:qtable}) reveal a quality-aware learned
policy:
\begin{itemize}\itemsep0pt
  \item \textbf{Poor-quality regions} (gap $>$30\%): relocate and swap dominate,
        effective at escaping large local optima.
  \item \textbf{High-quality regions} (gap $<$15\%): 2-opt dominates, consistent
        with its strength in fine-tuning near-optimal solutions.
  \item \textbf{Restart} is uniformly low across all bands --- the agent learns
        to treat it as a last resort.
\end{itemize}
This gradient from perturbation-preference to refinement-preference is exactly
what an expert practitioner would engineer manually, here discovered autonomously.

\begin{figure}[!htbp]
  \centering
  \begin{subfigure}{0.48\linewidth}
    \includegraphics[width=\linewidth]{tour_eil51.png}
    \caption{eil51}
  \end{subfigure}
  \hfill
  \begin{subfigure}{0.48\linewidth}
    \includegraphics[width=\linewidth]{tour_berlin52.png}
    \caption{berlin52}
  \end{subfigure}
  \caption{Best tours found by Q-Learning. Routes are visually clean with
           minimal crossing edges.}
  \label{fig:tours}
\end{figure}

\subsection{Generalization Study and State Encoding Limitation}

\begin{table}[!htbp]
\centering
\caption{Generalization Test: Mean Gap (\%) on unseen instances.}
\label{tab:gen}
\small
\begin{tabular}{lcc}
\toprule
\textbf{Test Instance} & \textbf{Pre-trained} & \textbf{Fresh Q-Learning} \\
\midrule
st70    & 4.06$\pm$1.48 & \textbf{3.87}$\pm$1.18 \\
kroA100 & 8.17$\pm$1.58 & \textbf{6.66}$\pm$1.69 \\
\bottomrule
\end{tabular}
\end{table}

\begin{figure}[!htbp]
  \centering
  \includegraphics[width=0.78\linewidth]{generalization_bar.png}
  \caption{Pre-trained (eil51+berlin52) vs.\ Fresh Q-Learning on test instances.
           Fresh training consistently outperforms direct Q-table transfer.}
  \label{fig:genbar}
\end{figure}

As Table~\ref{tab:gen} and Figure~\ref{fig:genbar} show, pre-training on
smaller instances does \emph{not} improve performance on larger unseen ones.
This negative transfer highlights a critical limitation in our MDP formulation: 
the \emph{quality\_band} encoding functions as an "oracle" feature, directly 
relying on each instance's known optimal TSPLIB value. When transferred to a new 
instance landscape, this static, instance-dependent reference causes a state distribution shift.

%--------------------------------------------------------------------
\section{Conclusion}

In this project, we presented a tabular Q-learning based hyper-heuristic for
adaptive operator selection in TSP local search. Based on the experimental
results, the main observations can be summarized as follows:
\begin{enumerate}\itemsep0pt
\item The adaptive Q-learning approach generally performs better than the
fixed-operator baselines by choosing different operators during the
search process.
\item The adaptive stagnation-based epsilon strategy gives better results
than the other exploration strategies tested in this study.
\item The learned operator transitions show that some operator sequences,
such as \emph{swap}$\to$\emph{2-opt}, can be useful during the search.
\item The Q-table shows that the agent changes its operator choices depending on the quality of the current solution.
\item The transfer learning results suggest that normalized state features
may be necessary for better generalization across different TSP
instances.
\end{enumerate}

As future work, the state representation can be improved by using more
general relative measures instead of instance-specific values. In addition,
deep Q-network (DQN) based approaches can be tested for continuous state
features and larger problem instances.

%--------------------------------------------------------------------
\bibliographystyle{ieeetr}
\begin{thebibliography}{15}

\bibitem{garey79}
M.~R. Garey and D.~S. Johnson,
\textit{Computers and Intractability: A Guide to the Theory of NP-Completeness}.
Freeman, 1979.

\bibitem{applegate06}
D.~Applegate, R.~Bixby, V.~Chvatal, and W.~Cook,
\textit{The Traveling Salesman Problem: A Computational Study}.
Princeton University Press, 2006.

\bibitem{helsgott00}
K.~Helsgott, O.~Rooij, and R.~Wiegand,
``LKH: An Effective Implementation of the Lin-Kernighan Heuristic,''
\textit{European Journal of Operational Research}, vol.~126, no.~1, pp.~106--130, 2000.

\bibitem{burke10}
E.~K. Burke \textit{et al.},
``A Classification of Hyper-heuristic Approaches,''
in \textit{Handbook of Metaheuristics}, Springer, 2010, pp.~449--468.

\bibitem{lourenco03}
H.~R. Lourenco, O.~C. Martin, and T.~Stutzle,
``Iterated Local Search,''
in \textit{Handbook of Metaheuristics}, Springer, 2003, pp.~320--353.

\bibitem{watkins92}
C.~J. C.~H. Watkins and P.~Dayan,
``Q-learning,''
\textit{Machine Learning}, vol.~8, no.~3--4, pp.~279--292, 1992.

\bibitem{meignan15}
D.~Meignan \textit{et al.},
``A Review and Taxonomy of Interactive Optimization Methods in Operations Research,''
\textit{ACM TIIS}, vol.~5, no.~3, pp.~1--43, 2015.

\bibitem{dacosta08}
L.~Da Costa, F.~Guimaraes, and A.~Delbem,
``Operator Selection with Dynamic Multi-Armed Bandits,''
in \textit{Proc. GECCO 2008}, ACM, pp.~1539--1546.

\bibitem{lin65}
S.~Lin,
``Computer Solutions of the Traveling Salesman Problem,''
\textit{Bell System Technical Journal}, vol.~44, no.~10, pp.~2245--2269, 1965.

\bibitem{lin73}
S.~Lin and B.~W. Kernighan,
``An Effective Heuristic Algorithm for the Traveling Salesman Problem,''
\textit{Operations Research}, vol.~21, no.~2, pp.~498--516, 1973.

\bibitem{or76}
I.~Or,
\textit{Traveling Salesman-type Combinatorial Problems and Their Relation to the
Logistics of Regional Blood Banking}.
Ph.D. Dissertation, Northwestern University, 1976.

\bibitem{bello16}
I.~Bello \textit{et al.},
``Neural Combinatorial Optimization with Reinforcement Learning,''
\textit{arXiv:1611.09940}, 2016.

\bibitem{kool19}
W.~Kool, H.~van Hoof, and M.~Welling,
``Attention, Learn to Solve Routing Problems!''
in \textit{Proc. ICLR 2019}.

\bibitem{hottung20}
A.~Hottung and K.~Tierney,
``Neural Large Neighbourhood Search for the Capacitated Vehicle Routing Problem,''
in \textit{Proc. ECAI 2020}, pp.~443--450.

\bibitem{tokic10}
M.~Tokic,
``Adaptive $\varepsilon$-greedy Exploration in Reinforcement Learning Based on
Value Differences,''
in \textit{Proc. KI 2010}, Springer, pp.~203--210.

\end{thebibliography}

\end{document}
